\section{Introduction}\label{sec:introduction}

% An introduction where you explain the aims and rationale for the physics case
% and what you have done. At the end of the introduction you should give a brief
% summary of the structure of the report.
%
% State explicitly what is your own code and what is scikit-learn/other
% libraries. End with a paragraph on the report's structure and the repo link.
% No results in the introduction.

In the essay \citetitle{sutton_2019}, Richard \textcite{sutton_2019}
argues that the most significant advances in artificial intelligence
have ultimately come from general methods that leverage computation,
rather than approaches relying on extensive human-engineered domain
knowledge. This observation highlights a central idea in modern machine
learning: the ability to learn patterns and representations directly
from data, rather than explicitly specifying how a problem should be
solved. Artificial neural networks are one such approach, providing a
flexible framework for learning complex relationships from data.

Linear and logistic regression provide relatively simple approaches
to regression and classification, respectively. While these models
offer interpretable and computationally efficient solutions, their
predictive capabilities are constrained by their functional form.
Neural networks extend this framework through compositions of
parameterized transformations and nonlinear activation functions,
allowing them to represent more complex relationships between inputs
and outputs. This increased flexibility, however, also introduces
challenges related to model architecture, optimization, and
generalization.

The objective of this project is to implement and investigate
feed-forward neural networks (FFNNs) for regression and classification.
In particular, we study the effects of network architecture,
activation functions, and optimization algorithms on model performance
and training behavior. We also compare the resulting models with
linear and logistic regression approaches,
examining the trade-offs between model complexity, predictive
performance, and computational cost.

The remainder of this report is structured as follows.
\Cref{sec:method} presents the theoretical background and methods
used to construct and train the networks.
\Cref{sec:implementation} describes the datasets, preprocessing,
implementation details, and verification of the code.
\Cref{sec:results-discussion} presents and discusses the results,
including comparisons with the linear and logistic regression models.
Finally, \Cref{sec:conclusion} summarizes the main findings and
outlines possible directions for future work. Appendix
\ref{app:ai-tools} documents our use of large language models.
The code and results underlying the figures and tables are available
in our publicly hosted repository~\cite{riseth_2026a}.
